% ============================================================
% PRN232 - Section 8: Deployment (Vietnamese)
% ============================================================
\section{\bi{Triển khai và tích hợp liên tục}{Deployment and Continuous Integration}}
\label{sec:deploy}

\subsection{\bi{Kiến trúc triển khai}{Deployment Architecture}}

\bi{
Hệ thống được triển khai trên Vercel dưới dạng ứng dụng serverless, tách
thành ba dự án độc lập: back-end NestJS, cổng web Next.js và một dịch
vụ cố định. Ba URL sản xuất đã được kiểm tra trả về HTTP 200.
}{
The system is deployed on Vercel as a serverless setup, split into
three independent projects: the NestJS back end, the Next.js portal
and a static service. All three production URLs have been verified to
return HTTP 200.
}

\begin{table}[t]
\caption{\bi{Các điểm cuối sản xuất và trạng thái kiểm tra ngày 28/10/2026.}{Production endpoints and verification status as of 28 October 2026.}}
\label{tab:deploy}
\centering
\footnotesize
\begin{tabular}{@{}lll@{}}
\toprule
\textbf{\bi{Dịch vụ}{Service}} & \textbf{URL} & \textbf{\bi{Trạng thái}{Status}} \\
\midrule
\texttt{omnicast.vercel.app}      & \texttt{https://omnicast.vercel.app}      & HTTP 200 \\
\texttt{omnicast-fe.vercel.app}    & \texttt{https://omnicast-fe.vercel.app}    & HTTP 200 \\
\texttt{omnicast-api.vercel.app}   & \texttt{https://omnicast-api.vercel.app}   & HTTP 200 \\
\bottomrule
\end{tabular}
\end{table}

\subsection{\bi{Quy trình CI/CD}{The CI/CD Pipeline}}

\bi{
Tệp \texttt{.github/workflows/ci.yml} định nghĩa năm công việc chạy
song song và một cổng chất lượng. Ba công việc kiểm thử chạy độc lập
trên nền \texttt{ubuntu-latest} với Node.js 20, Java 17 và Flutter
ổn định. Công việc tổng hợp gom artefact của cả ba, sinh báo cáo HTML
tương tác, và dừng pipeline nếu bất kỳ bộ kiểm thử nào thất bại.
}{
The file \texttt{.github/workflows/ci.yml} defines five jobs running
in parallel plus a quality gate. Three test jobs run independently on
\texttt{ubuntu-latest} with Node.js 20, Java 17 and stable Flutter. The
aggregation job downloads all three artefact sets, generates the
interactive HTML report, and fails the pipeline if any suite fails.
}

\begin{figure*}[t]
\centering
\fitfull{\figCicd}
\caption{\bi{Quy trình CI/CD năm công việc. Ba công việc kiểm thử chạy
song song, hợp nhất tại công việc tổng hợp và chỉ triển khai khi cổng
chất lượng đạt.}{The five-job CI/CD pipeline. Three test jobs run in
parallel, merge at the summary job, and deployment only proceeds once
the quality gate passes.}}
\label{fig:cicd}
\end{figure*}

\bi{
Công việc triển khai chỉ chạy khi nhánh là \texttt{main} và sự kiện
là \texttt{push}; nó dùng \texttt{vercel-action} để đẩy back-end và
cổng web, truyền \texttt{NEXT\_PUBLIC\_API\_URL} từ secret của kho.
Nhờ đặt \texttt{continue-on-error} cho công việc triển khai, một sự cố
triển khai không xoá artefact kiểm thử đã tạo.
}{
The deploy job runs only on the \texttt{main} branch on a
\texttt{push} event; it uses \texttt{vercel-action} to push the back
end and the portal, passing \texttt{NEXT\_PUBLIC\_API\_URL} from a
repository secret. Because the deploy job sets
\texttt{continue-on-error}, a deployment failure does not destroy the
test artefacts already produced.
}

\subsection{\bi{Cấu hình biến môi trường}{Environment Configuration}}

\bi{
Mẫu biến môi trường nằm ở \texttt{backend/.env.example} và gồm
bốn nhóm: kết nối cơ sở dữ liệu Supabase kèm mật khẩu cần thay thế,
khóa \texttt{JWT\_SECRET} cùng thời hạn token, các khóa API bên thứ
ba (OpenAI, TMDB, YouTube, Twitch, Cloudinary), và cấu hình MediaMTX
cho môi trường cục bộ. Tệp \texttt{.env} thực sự không bao giờ được
đưa vào kiểm soát phiên bản.
}{
The environment template lives at \texttt{backend/.env.example} and
covers four groups: the Supabase database connection with a password
placeholder, the \texttt{JWT\_SECRET} and token lifetimes, the
third-party API keys (OpenAI, TMDB, YouTube, Twitch, Cloudinary), and
the local MediaMTX configuration. The real \texttt{.env} file is
never committed.
}

\subsection{\bi{Bảo mật triển khai}{Deployment Security}}

\bi{
Tám biện pháp bảo mật được áp dụng ở tầng ứng dụng: \texttt{Helmet}
thiết lập các tiêu đề bảo mật HTTP; \texttt{ThrottlerModule} giới hạn
100 yêu cầu mỗi phút; CORS dùng danh sách cho phép tường minh kèm
kiểm tra miền Vercel và mạng nội bộ; \texttt{ValidationPipe} loại bỏ
trường không khai báo; \texttt{HttpExceptionFilter} chuẩn hoá lỗi;
mật khẩu được băm bằng bcrypt 12 vòng; refresh token được xoay vòng
và thu hồi được lưu trong cơ sở dữ liệu; và phương thức bất đồng bộ
được bọc trong \texttt{promise} không chặn để tránh lỗi nghiệp vụ do
lỗi hạ tầng ghi log.
}{
Eight application-layer security measures are in place: \texttt{Helmet}
sets the HTTP security headers; \texttt{ThrottlerModule} caps traffic at
100 requests per minute; CORS uses an explicit allow-list together
with Vercel-domain and private-network checks; \texttt{ValidationPipe}
strips undeclared fields; \texttt{HttpExceptionFilter} normalises
errors; passwords are hashed with 12-round bcrypt; refresh tokens are
rotated and revocations are persisted; and the asynchronous audit
write is wrapped in a non-awaited promise so an infrastructure
failure cannot break a business operation.
}

\subsection{\bi{Kiểm thử tải và độ bền}{Load and Resilience Testing}}

\bi{
Bộ kiểm thử tải k6 được sử dụng để kiểm tra hành vi dưới tải đồng
thời của nhiều người xem, tập trung vào các endpoint EPG vì đây là
đường đi nóng nhất của hệ thống. Kết quả cho thấy các endpoint đọc
giữ được độ trễ ổn định trong khi endpoint ghi vẫn đáp ứng thời gian
ghi nhật ký kiểm toán ở mức không chặn luồng.
}{
A k6 load-test suite exercises the system under concurrent viewers,
focusing on the EPG endpoints because they are the hottest path.
The results show that read endpoints maintain stable latency while
write endpoints still meet their response-time target thanks to the
non-blocking audit write.
}
